\documentclass[11pt,a4paper]{article}

% ============================================================
% Relatório preliminar — NeuraTrade
% Detecção de anomalias em séries financeiras da B3 com LSTM-Autoencoder
%
% Documento VIVO: a cada fechamento de milestone, acrescenta-se uma seção
% em sections/ e registra-se aqui via \input. Servirá de base para o
% artigo final.
%
% Compilar:  tectonic main.tex      (ou latexmk -pdf main.tex)
% ============================================================

\usepackage[utf8]{inputenc}
\usepackage[T1]{fontenc}
\usepackage[brazil]{babel}
\usepackage[a4paper,margin=2.5cm]{geometry}
\usepackage{amsmath,amssymb}
\usepackage{booktabs}
\usepackage{graphicx}
\usepackage{hyperref}
\usepackage{xurl}
\usepackage{enumitem}
\usepackage{authblk}
\usepackage{caption}
\usepackage[numbers,sort&compress]{natbib}

\hypersetup{colorlinks=true,linkcolor=blue,citecolor=blue,urlcolor=blue}

\title{\textbf{Detecção Não Supervisionada de Anomalias em Séries Temporais
Financeiras Brasileiras com LSTM-Autoencoder}\\[0.3em]
\large Relatório Preliminar}

\author[ ]{Ana Beatriz}
\author[ ]{Miguel Cerne}
\affil[ ]{Redes Neurais (2026.1)}

\date{\today}

\begin{document}
\maketitle

\begin{abstract}
A detecção de anomalias em mercados financeiros é dominada, na literatura, pelo
mercado norte-americano e por criptomoedas. Este trabalho aplica uma arquitetura
\emph{LSTM-Autoencoder} à detecção não supervisionada de anomalias em quatro ações
da B3 (PETR4, VALE3, AMER3 e ITUB4), cobrindo distintos setores econômicos. O modelo
é treinado apenas sobre o período considerado de ``normalidade'' (2010--2019) e sinaliza
como anômalas as janelas cujo erro de reconstrução excede um limiar; comparam-se
limiares estático (percentil fixo) e dinâmico (percentil em janela móvel causal).
A avaliação combina validação narrativa (sobreposição com eventos econômicos e políticos
brasileiros documentados) e quantitativa (injeção controlada de anomalias sintéticas,
com Precision, Recall e F1). Este relatório preliminar documenta, de forma incremental e
academicamente rastreável, as decisões metodológicas, suas fontes e as observações de
cada etapa do projeto. Até o momento concluíram-se as etapas de \emph{infraestrutura}
(M0) e de \emph{coleta e análise exploratória} (M1).

\end{abstract}

\noindent\textbf{Palavras-chave:} detecção de anomalias; LSTM-Autoencoder;
séries temporais financeiras; B3; aprendizado não supervisionado.

\tableofcontents

\section{Introdução}

A identificação de comportamentos atípicos em séries de preços de ativos é relevante
tanto para a gestão de risco quanto para a compreensão de como eventos econômicos e
políticos se refletem nos mercados. A maior parte da literatura de detecção de anomalias
com aprendizado profundo concentra-se no mercado norte-americano e em
criptomoedas~\cite{li2020lstm,bitcoin2021anomaly}, deixando sub-representado o contexto
de mercados emergentes como o brasileiro.

Este trabalho investiga a detecção \emph{não supervisionada} de anomalias em ações da
B3 por meio de um \emph{LSTM-Autoencoder}~\cite{hochreiter1997lstm}. A premissa é que um
autoencoder treinado apenas em períodos de comportamento típico aprende a reconstruir os
padrões normais da série; diante de movimentos atípicos, o erro de reconstrução cresce e
sinaliza a anomalia.

\subsection{Contribuições}
O projeto combina quatro contribuições:
\begin{enumerate}[noitemsep]
  \item Aplicação a ativos nacionais (B3), contexto sub-representado na literatura.
  \item Correlação sistemática e documentada entre anomalias detectadas e eventos
        históricos brasileiros.
  \item Comparação entre limiar \emph{estático} (percentil fixo) e \emph{dinâmico}
        (percentil em janela móvel).
  \item Avaliação de generalização entre setores econômicos distintos.
\end{enumerate}

\subsection{Ativos e período}
Foram selecionados quatro ativos, um por setor, conforme a Tabela~\ref{tab:ativos}.
O período de coleta vai de 2010 a 2024; o treino (``normalidade'') usa 2010--2019 e o
teste, 2020--2024.

\begin{table}[h]
\centering
\caption{Ativos analisados.}
\label{tab:ativos}
\begin{tabular}{lll}
\toprule
\textbf{Ticker} & \textbf{Empresa} & \textbf{Setor} \\
\midrule
PETR4.SA & Petrobras PN    & Energia/commodities \\
VALE3.SA & Vale S.A.       & Mineração \\
AMER3.SA & Americanas S.A. & Varejo \\
ITUB4.SA & Itaú Unibanco   & Financeiro \\
\bottomrule
\end{tabular}
\end{table}

\subsection{Sobre este documento}
Trata-se de um relatório preliminar e \emph{vivo}: a cada fechamento de \emph{milestone}
acrescenta-se uma seção descrevendo, academicamente, as decisões tomadas, suas fontes,
as observações empíricas e as conclusões parciais. As decisões de configuração são
registradas em paralelo como \emph{Architecture Decision Records} (ADRs) no
repositório (\texttt{docs/adr/}), referenciados ao longo do texto.

\section{Trabalhos Relacionados}

A detecção de anomalias com autoencoders recorrentes baseia-se na ideia de que a
reconstrução de padrões aprendidos falha sobre observações fora da distribuição de
treino. \citet{li2020lstm}, em referência amplamente reproduzida, aplica um
LSTM-Autoencoder a preços de ações (Johnson \& Johnson), adotando janelas de 30 passos
de tempo e definindo o limiar de anomalia a partir do erro de reconstrução do treino.
\citet{valkov2019lstm} fornece uma implementação canônica em Keras, com escolhas de
projeto que servem de referência prática (camadas LSTM de 64 unidades, \emph{dropout}
de 0{,}2, \texttt{validation\_split} de 0{,}1 e, crucialmente, \texttt{shuffle=False}
para preservar a ordem temporal). \citet{petrovic2019lstm} disponibiliza um repositório
com abordagem equivalente.

Trabalhos mais recentes ampliam o desenho experimental. \citet{liu2025robust} propõem um
arcabouço híbrido que acopla um LSTM-Autoencoder a redes adversárias generativas e
introduzem explicitamente um \emph{mecanismo de injeção artificial de anomalias} para
avaliação quantitativa, validado em múltiplas categorias de ações --- desenho que
inspira diretamente a avaliação setorial e por injeção adotada neste projeto. No domínio
de criptomoedas, \citet{bitcoin2021anomaly} exploram a mesma família de modelos sobre
valores do Bitcoin.

Em relação a esses trabalhos, a principal lacuna endereçada aqui é a aplicação ao mercado
brasileiro (B3) com correlação documentada a eventos nacionais e comparação explícita
entre estratégias de limiar estático e dinâmico.

% \citet usado acima requer natbib; caso o estilo não suporte, ver nota em report/README.

\section{Metodologia}
\label{sec:metodologia}

Esta seção descreve o desenho metodológico do projeto. Os valores numéricos de
hiperparâmetros e suas justificativas são apresentados nas seções de cada
\emph{milestone} e consolidados nos ADRs do repositório.

\subsection{Pré-processamento}
A feature principal é o \emph{retorno logarítmico diário},
$r_t = \ln(P_t / P_{t-1})$, sobre o preço de fechamento ajustado. Adota-se um
\textbf{split temporal antes da normalização}: o \texttt{MinMaxScaler} é ajustado
\emph{apenas} sobre o treino e então aplicado ao teste, evitando vazamento de
informação do futuro. A série é então segmentada em \emph{janelas deslizantes} de
tamanho fixo (ver Seção~\ref{sec:m1}).

\subsection{Modelo}
A arquitetura é um \emph{LSTM-Autoencoder}~\cite{hochreiter1997lstm}: um codificador
LSTM comprime a janela em uma representação latente (gargalo), e um decodificador LSTM
reconstrói a sequência, otimizando o erro quadrático médio (MSE) de reconstrução. Treina-se
\textbf{um modelo por ativo}, condição necessária para a comparação setorial.

\subsection{Detecção}
Uma janela é marcada como anômala quando seu erro de reconstrução ultrapassa um limiar.
Comparam-se duas estratégias:
\begin{itemize}[noitemsep]
  \item \textbf{Estático:} percentil fixo (95) do erro de reconstrução do treino.
  \item \textbf{Dinâmico:} percentil em janela móvel \emph{causal} (apenas o passado),
        sensível a mudanças de regime.
\end{itemize}

\subsection{Avaliação}
Por ser um problema não supervisionado, a avaliação combina duas vias~\cite{liu2025robust}:
\begin{itemize}[noitemsep]
  \item \textbf{Narrativa:} verificação de que as anomalias detectadas correspondem a
        eventos reais conhecidos (e.g., \emph{crash} de março/2020; caso Americanas em
        janeiro/2023).
  \item \textbf{Quantitativa:} injeção controlada de anomalias sintéticas (\emph{price
        shocks} e \emph{volume spikes}) em posições conhecidas, permitindo calcular
        Precision, Recall e F1.
\end{itemize}

\paragraph{Nota metodológica.} O período de treino (2010--2019) contém eventos de forte
volatilidade (recessão de 2014--2016, Operação Lava Jato, \emph{impeachment} de 2016).
A definição de ``normalidade'' é, portanto, relativa, e essa limitação é assumida
explicitamente.

\section{M0 --- Infraestrutura e Reprodutibilidade}
\label{sec:m0}

A primeira etapa estabeleceu a infraestrutura que garante reprodutibilidade e
separação entre lógica e orquestração: a lógica reside em \texttt{src/} e os notebooks
apenas orquestram e visualizam.

\subsection{Ambiente e dependências}
Fixou-se o interpretador em \textbf{Python 3.12.13} (via \texttt{pyenv}), versão idêntica
à do Google Colab à época, alvo principal de execução. As dependências foram
\emph{pinadas} (TensorFlow/Keras~\cite{tensorflow2015}, \texttt{yfinance}, \texttt{pandas},
\texttt{numpy}, \texttt{scikit-learn}, \texttt{matplotlib}, \texttt{seaborn}), habilitando
instalação editável (\texttt{pip install -e .}). Notebooks têm suas saídas removidas
automaticamente no versionamento (\texttt{nbstripout}), reduzindo ruído em \emph{diffs}.

\subsection{Configuração centralizada e seeds}
Todos os hiperparâmetros residem em \texttt{config.yaml}, fonte única de verdade: alterar
um experimento significa alterar esse arquivo, não os notebooks. O módulo
\texttt{src/config.py} carrega essa configuração e expõe \texttt{set\_seeds()}, que fixa
as sementes de \texttt{random}, \texttt{numpy} e \texttt{tensorflow} para reprodutibilidade.

\subsection{Decisões de hiperparâmetros e proveniência}
Um princípio metodológico central do projeto é a \textbf{rastreabilidade da proveniência}
de cada hiperparâmetro. Cada escolha é classificada como: \emph{fundamentada}
(ancorada no enunciado do projeto ou na literatura), \emph{default de literatura}
(convenção da implementação de referência), \emph{decisão de projeto} (escolha de
\emph{design} sem fonte direta, com \emph{rationale} explícito) ou \emph{provisória}
(valor a ser calibrado por experimento). A Tabela~\ref{tab:hparams} consolida as
principais decisões; os valores provisórios não devem sustentar resultados finais sem
validação.

\begin{table}[h]
\centering
\caption{Principais hiperparâmetros e sua proveniência (consolidado dos ADRs).}
\label{tab:hparams}
\small
\begin{tabular}{lll}
\toprule
\textbf{Parâmetro} & \textbf{Valor} & \textbf{Proveniência} \\
\midrule
\texttt{window\_size}        & 30      & Fundamentada~\cite{li2020lstm,valkov2019lstm} \\
\texttt{lstm\_units}         & 64      & Default de literatura~\cite{valkov2019lstm} \\
\texttt{dropout}             & 0{,}2   & Default de literatura~\cite{valkov2019lstm} \\
\texttt{latent\_dim}         & 16      & Decisão de projeto (calibrar) \\
\texttt{loss}                & MSE     & Fundamentada~\cite{li2020lstm} \\
\texttt{shuffle}             & False   & Fundamentada (causalidade)~\cite{valkov2019lstm} \\
\texttt{validation\_split}   & 0{,}1   & Default de literatura~\cite{valkov2019lstm} \\
\texttt{threshold\_percentile} & 95    & Fundamentada (enunciado) \\
\texttt{dynamic\_window}     & 252     & Provisória ($\approx$1 ano útil) \\
\texttt{n\_injections}       & 50      & Provisória \\
\bottomrule
\end{tabular}
\end{table}

\subsection{Riscos metodológicos registrados}
Três riscos foram documentados já nesta fase, para serem honrados na implementação:
(i) o \texttt{validation\_split} do Keras separa o \emph{bloco final} dos dados antes de
embaralhar --- logo \texttt{shuffle=False} é obrigatório, sob pena de vazamento temporal;
(ii) o limiar dinâmico deve ser \emph{causal}, usando apenas observações passadas;
(iii) o \texttt{MinMaxScaler} ajustado no treino faz com que extremos do teste possam
exceder o intervalo $[0,1]$ --- comportamento esperado e desejável, pois é onde residem as
anomalias, mas que torna o limiar sensível a \emph{outliers} isolados (motivando o uso de
percentil em vez do máximo).

\section{M1 --- Coleta e Análise Exploratória}
\label{sec:m1}

\subsection{Coleta e cache}
Os dados são obtidos via \texttt{yfinance} com \texttt{auto\_adjust=True}, de modo que o
preço de fechamento já incorpora ajustes por \emph{splits} e dividendos. A arquitetura de
dados separa estritamente o acesso à rede da leitura: as funções \texttt{cache\_*} baixam
e persistem os dados em \texttt{data/raw/} (CSV, versionado), enquanto \texttt{load\_*}
--- usadas por todo o \emph{pipeline} --- leem apenas do cache, jamais da API. Isso garante
execução \emph{offline} e reprodutível dos notebooks.

\subsection{Integridade das séries}
As quatro séries apresentam \textbf{3724 pregões} cada (2010-01-04 a 2024-12-30), sem
valores ausentes (NaN) após o ajuste. A Tabela~\ref{tab:integridade} resume os achados.
A função \texttt{integrity\_report} computa, por ativo, contagens de NaN, dias com volume
zero (halts/iliquidez) e os extremos de log-retorno diário, que funcionam como
indicadores tanto de anomalias reais quanto de eventuais artefatos.

\begin{table}[h]
\centering
\caption{Integridade e maior queda diária de log-retorno por ativo.}
\label{tab:integridade}
\small
\begin{tabular}{lcccl}
\toprule
\textbf{Ticker} & \textbf{Pregões} & \textbf{NaN} & \textbf{Vol.\ zero} & \textbf{Maior queda (data)} \\
\midrule
PETR4.SA & 3724 & 0 & 54 & $-0{,}352$ (2020-03-09) \\
VALE3.SA & 3724 & 0 & 27 & $-0{,}282$ (2019-01-28) \\
AMER3.SA & 3724 & 0 & 25 & $-1{,}484$ (2023-01-12) \\
ITUB4.SA & 3724 & 0 & 34 & $-0{,}198$ (2021-10-04) \\
\bottomrule
\end{tabular}
\end{table}

\subsection{Anomalias reais identificadas}
A EDA já revela eventos extremos que servirão de âncora para a validação narrativa (M6):
\begin{itemize}[noitemsep]
  \item \textbf{PETR4}, 09/03/2020: queda associada ao início da pandemia de COVID-19 e à
        guerra de preços do petróleo.
  \item \textbf{VALE3}, 28/01/2019: reação ao rompimento da barragem de Brumadinho.
  \item \textbf{AMER3}, 12/01/2023: divulgação da fraude contábil da Americanas (queda de
        $\approx$78\% no dia).
  \item \textbf{ITUB4}, 04/10/2021: maior queda do setor financeiro no período.
\end{itemize}
A distribuição dos log-retornos exibe caudas pesadas (leptocurtose), o que reforça a
escolha de limiares por percentil em vez de critérios baseados em desvio-padrão.

\subsection{O caso AMER3}
A série da Americanas merecia atenção especial por causa da reorganização societária e do
colapso pós-fraude. Dois pontos foram esclarecidos pela EDA:
\begin{enumerate}[noitemsep]
  \item \textbf{Histórico de treino existe.} Diferentemente da hipótese inicial de ausência
        de dados em 2010--2019, o \texttt{yfinance} fornece a série desde 2010 por
        continuidade de \emph{ticker} (herdada da Lojas Americanas). Há, portanto, dados de
        ``normalidade'' para o ativo.
  \item \textbf{Sinal vs.\ artefato.} A queda de 12/01/2023 é uma anomalia \emph{real} e
        deve ser preservada (é justamente o evento-alvo). Já o salto de log-retorno
        $\approx +1{,}03$ em 14/11/2024 é compatível com um \emph{grupamento} (\emph{reverse
        split}) não totalmente ajustado pela fonte, configurando provável \textbf{artefato}.
\end{enumerate}
Decidiu-se, para a etapa de pré-processamento (M2), \textbf{corrigir o fator de grupamento}
do AMER3, em vez de recortar o período --- preservando a continuidade temporal e o caso
completo do ativo.

\section{M2 --- Pré-processamento}
\label{sec:m2}

Esta etapa transforma as séries de preço na representação consumida pelo modelo,
respeitando a ordem metodológica que evita vazamento temporal (Seção~\ref{sec:metodologia}).

\subsection{Log-retornos}
A feature principal é o retorno logarítmico diário $r_t = \ln(P_t/P_{t-1})$, calculado
sobre o fechamento ajustado. Por ser invariante a transformações multiplicativas de
escala, é imune a fatores de \emph{split}/grupamento já aplicados ao preço --- propriedade
relevante para o caso AMER3 (Seção~\ref{sec:m2-amer3}).

\subsection{Split temporal e normalização}
Aplica-se primeiro o \emph{split} temporal (treino até 2019; teste a partir de 2020) e só
então a normalização: o \texttt{MinMaxScaler} é ajustado \emph{exclusivamente} no treino e
aplicado a ambas as partições. Como consequência esperada, uma fração dos pontos de teste
extrapola o intervalo $[0,1]$ --- precisamente os extremos pós-2020 fora da faixa observada
no treino. Esse comportamento é desejável: é onde se concentram as anomalias, e o erro de
reconstrução tende a crescer ali. As janelas deslizantes ($\texttt{window\_size}=30$) são
geradas \emph{dentro} de cada partição, nunca cruzando a fronteira treino/teste. O
resultado são $\approx 2450$ janelas de treino e $\approx 1215$ de teste por ativo, no
formato $(n, 30, 1)$.

\subsection{Revisão da hipótese sobre o AMER3}
\label{sec:m2-amer3}
Em M1, levantou-se a hipótese de que o salto de log-retorno de $\approx +1{,}03$ em
14/11/2024 seria um \emph{artefato} de grupamento, e decidiu-se, provisoriamente,
``corrigir o fator''. A investigação conduzida em M2 \textbf{refutou essa hipótese} e
ilustra a importância de validar premissas com os dados:
\begin{itemize}[noitemsep]
  \item Os grupamentos reais do AMER3 ocorreram em \textbf{julho e agosto de 2024} (quatro
        eventos de razão 1:100 registrados pela fonte) e \textbf{já são corrigidos por}
        \texttt{auto\_adjust=True}: o log-retorno ajustado nessas datas é pequeno
        ($+0{,}09$; $+0{,}01$; $-0{,}18$; $+0{,}34$), com série contínua.
  \item Em 14/11/2024 o preço bruto saltou de R\$\,3{,}36 para R\$\,9{,}41 com volume
        cerca de três vezes a média e \emph{sem} \emph{split} registrado --- caracterizando
        um \textbf{movimento de mercado real} (especulação durante a recuperação judicial),
        assim como o salto de 15/08/2024.
  \item Como o log-retorno é invariante a escala, \emph{splits} já aplicados não produzem
        retorno espúrio.
\end{itemize}
\textbf{Decisão revista:} nenhuma correção é aplicada; o AMER3 segue o mesmo pipeline dos
demais ativos, e seus saltos de 2024 são tratados como anomalias reais (verdade narrativa
para M6). O registro completo está no ADR-0007.

\section{M3 --- Modelo e Treino}
\label{sec:m3}

\subsection{Arquitetura}
Implementou-se o LSTM-Autoencoder descrito na Seção~\ref{sec:metodologia}: um codificador
\texttt{LSTM(64)} seguido de \emph{dropout} e de uma camada densa de gargalo
(\texttt{latent\_dim=16}); o decodificador replica o vetor latente (\texttt{RepeatVector})
e o reconstrói com \texttt{LSTM(64)} e \texttt{TimeDistributed(Dense(1))}. A rede tem
$\approx$38{,}7 mil parâmetros --- deliberadamente pequena, coerente com o tamanho do
conjunto ($\approx$2450 janelas por ativo). A perda é o MSE de reconstrução e o otimizador
é o Adam.

\subsection{Protocolo de treino}
Treinou-se \textbf{um modelo por ativo}, com validação por bloco cronológico
(\texttt{shuffle=False}) e \texttt{EarlyStopping} (paciência 10, \texttt{restore\_best\_weights}).
Os pesos são salvos em \texttt{models/<ticker>.keras} (regeneráveis, não versionados). A
Tabela~\ref{tab:treino} resume a convergência: a parada antecipada atuou em todos os casos,
com perdas de validação da mesma ordem de grandeza entre ativos e sem sinais de
sobreajuste nas curvas.

\begin{table}[h]
\centering
\caption{Convergência do treino por ativo (\texttt{EarlyStopping}).}
\label{tab:treino}
\begin{tabular}{lcc}
\toprule
\textbf{Ticker} & \textbf{Épocas até parar} & \textbf{val\_loss final} \\
\midrule
PETR4.SA & 31 & 0{,}00352 \\
VALE3.SA & 15 & 0{,}00430 \\
AMER3.SA & 60 & 0{,}00354 \\
ITUB4.SA & 17 & 0{,}00440 \\
\bottomrule
\end{tabular}
\end{table}

\subsection{Sensibilidade do gargalo}
A dimensão latente havia sido fixada como decisão de projeto, a calibrar nesta etapa. Um
estudo em PETR4 com $\texttt{latent\_dim} \in \{8, 16, 32\}$ produziu \texttt{val\_loss}
praticamente idêntico (0{,}00342; 0{,}00343; 0{,}00342). Conclui-se que, neste regime, o
gargalo \textbf{não é restrição ativa} --- o modelo reconstrói os log-retornos igualmente
bem nas três configurações. Mantém-se \texttt{latent\_dim=16}, agora confirmado
empiricamente (registro no ADR-0003).

\subsection{Baseline de comparação}
Como chão de comparação sem aprendizado (issue \#14), implementou-se um \emph{z-score} em
janela móvel causal sobre os log-retornos. O baseline acende corretamente nas anomalias
reais --- por exemplo, no AMER3, $|z|=13{,}9$ em 12/01/2023 (fraude) e $|z|=9{,}7$ em
14/11/2024. Trata-se de um detector forte para eventos extremos; o valor agregado do
LSTM-Autoencoder será medido quantitativamente em M5, por injeção de anomalias sintéticas
mais sutis, onde se espera que heurísticas de cauda sejam menos eficazes.

\subsection{Erro de reconstrução}
Sobre o conjunto de treino (normalidade), o erro de reconstrução (MAE por janela) concentra-se
em valores baixos, com cauda à direita já indicando candidatos a anomalia. Essa distribuição
é a base para a definição formal dos limiares em M4.

\section{M4 --- Detecção e Limiares}
\label{sec:m4}

\subsection{Erro de reconstrução}
A detecção opera sobre o erro de reconstrução por janela --- o erro médio absoluto (MAE)
entre a janela de entrada e sua reconstrução. Optou-se pelo MAE por sua maior robustez a
\emph{outliers} isolados em relação ao MSE (este último disponível como opção). Uma janela
é classificada como anômala quando seu erro supera o limiar.

\subsection{Limiar estático}
O limiar estático é o percentil 95 do erro de reconstrução \emph{do treino}. Por ser
calculado sobre o período de normalidade (e não sobre o teste), não incorre em vazamento.
É simples e estável, mas, por definição, fixo: não se adapta a mudanças no regime de
volatilidade do ativo.

\subsection{Limiar dinâmico}
O limiar dinâmico é o percentil 95 calculado em uma janela móvel \textbf{causal} de
$\texttt{dynamic\_window}=252$ observações ($\approx$1 ano útil). A causalidade é garantida
excluindo-se o próprio ponto da janela (deslocamento de uma posição), de modo que uma
anomalia não infle o limiar que a avalia; os primeiros pontos, com janela incompleta, não
são marcados. Esse esquema constitui uma das contribuições metodológicas do trabalho.

\subsection{Comparação empírica}
A Tabela~\ref{tab:thresholds} compara, no conjunto de teste (2020--2024), a fração de
janelas marcadas por cada esquema. O resultado evidencia o comportamento esperado:
\begin{itemize}[noitemsep]
  \item Em ativos de regime relativamente estável (PETR4, VALE3), os dois limiares quase
        coincidem ($\approx$3\%).
  \item Sob mudança abrupta de regime, o limiar estático \textbf{satura}: no AMER3, todo o
        período pós-fraude/recuperação judicial faz o erro permanecer elevado, e o estático
        marca $\approx$27\% das janelas --- perdendo poder discriminativo. O limiar dinâmico,
        ao se recalibrar à volatilidade local, reduz a marcação para $\approx$14\% sem usar
        informação do futuro. Comportamento análogo no ITUB4 (11\% $\to$ 5\%).
\end{itemize}

\begin{table}[h]
\centering
\caption{Fração de janelas de teste marcadas como anômalas, por limiar.}
\label{tab:thresholds}
\begin{tabular}{lccc}
\toprule
\textbf{Ticker} & \textbf{Limiar estático} & \textbf{Estático (\%)} & \textbf{Dinâmico (\%)} \\
\midrule
PETR4.SA & 0{,}124 & 2{,}9 & 2{,}9 \\
VALE3.SA & 0{,}082 & 2{,}7 & 3{,}4 \\
AMER3.SA & 0{,}077 & 26{,}7 & 14{,}3 \\
ITUB4.SA & 0{,}088 & 11{,}0 & 4{,}9 \\
\bottomrule
\end{tabular}
\end{table}

\begin{figure}[h]
\centering
\includegraphics[width=\textwidth]{m4_erro_limiares}
\caption{Erro de reconstrução no teste (2020--2024) com limiar estático (p95, tracejado
vermelho) e dinâmico (verde). O dinâmico acompanha o regime de volatilidade; no AMER3 o
estático satura no período pós-fraude.}
\label{fig:m4-limiares}
\end{figure}

\paragraph{Observações metodológicas.} O percentil 95 implica, por construção, uma taxa-base
de $\approx$5\% de marcações --- a considerar na leitura das métricas (M5). Além disso, como
as janelas se sobrepõem, uma única anomalia gera uma rajada de detecções contíguas; antes de
confrontar com eventos reais (M6) ou calcular métricas (M5), detecções adjacentes devem ser
agrupadas em ``eventos''.

\subsection{Sensibilidade da janela dinâmica}
\label{sec:m4-sensibilidade}
O tamanho \texttt{dynamic\_window} era provisório; a Tabela~\ref{tab:sensibilidade} fecha a
questão variando-o em $\{126, 252, 504\}$. O comportamento é monotônico: a janela curta (126)
é reativa --- segue o ruído local e infla a marcação ($\approx$8\% nos ativos estáveis),
enquanto a janela longa (504) super-suaviza, aproximando-se de um limiar global e colapsando a
marcação a $\approx$0\% --- perdendo a capacidade de sinalizar. O valor intermediário
\textbf{252} produz taxas próximas da base-rate alvo ($\approx$5\% do p95) sem oscilar nem
colapsar, e é confirmado como escolha definitiva (ADR-0005).

\begin{table}[h]
\centering
\caption{Fração de janelas de teste marcadas pelo limiar dinâmico, por tamanho de janela.}
\label{tab:sensibilidade}
\begin{tabular}{lccc}
\toprule
\textbf{Ticker} & \textbf{126} & \textbf{252} & \textbf{504} \\
\midrule
PETR4.SA & 8{,}5\%  & 2{,}9\%  & 0{,}0\% \\
VALE3.SA & 8{,}2\%  & 3{,}4\%  & 0{,}3\% \\
AMER3.SA & 15{,}2\% & 14{,}3\% & 5{,}7\% \\
ITUB4.SA & 6{,}3\%  & 4{,}9\%  & 0{,}0\% \\
\bottomrule
\end{tabular}
\end{table}

\section{M5 --- Avaliação Quantitativa por Injeção Sintética}
\label{sec:m5}

\subsection{Motivação e desenho}
O problema é não supervisionado: não há rótulos de anomalia reais. Para obter
métricas quantitativas (\emph{Precision}, \emph{Recall}, F1) injetam-se anomalias em
posições \emph{conhecidas} da série de teste, produzindo um \emph{ground-truth}
controlado (ADR-0006). Foram inseridos $n=50$ \emph{price shocks} por ativo, em posições
aleatórias (\emph{seed} fixa), somando um delta de magnitude \texttt{price\_shock}$=0{,}15$
ao log-retorno \textbf{já escalado} --- o mesmo espaço em que o modelo foi treinado.
Injetar no preço bruto produziria choques de magnitude variável conforme o nível de preço
do ativo.

\subsection{Regra de rotulagem por janela}
Como o modelo processa janelas de 30 passos, um choque em qualquer posição eleva o erro de
reconstrução de toda a janela. Adota-se, portanto, a regra: uma janela é positiva se
\textbf{qualquer} passo injetado cair dentro dela. Os vetores de rótulo e de \emph{flags}
do detector ficam alinhados janela a janela, seguindo a mesma indexação de
\texttt{make\_windows} (Seção~\ref{sec:m2}), o que evita desalinhamento temporal no cálculo
das métricas.

\subsection{Resultados}
A Tabela~\ref{tab:metrics} apresenta as métricas por ativo, comparando limiar estático
(p95 do treino) e dinâmico (calculado sobre os erros da série perturbada). O padrão é
claro e metodologicamente informativo:
\begin{itemize}[noitemsep]
  \item O \emph{recall} é \textbf{uniformemente baixo} nos ativos de regime estável
        (PETR4, VALE3: recall $<7\%$). Com magnitude \emph{absoluta} de $0{,}15$ no espaço
        escalado, o choque sintético é comparável ao ruído normal desses ativos e raramente
        empurra a janela acima do limiar.
  \item Onde o regime de erro de base é mais alto (AMER3, ITUB4), os choques são detectados
        com mais frequência (AMER3 estático: recall $\approx$30\%, F1 $\approx$0{,}42).
  \item O limiar dinâmico tende a elevar a \emph{precision} (ex.\ VALE3 e ITUB4 acima de
        0{,}9), por se ajustar ao ruído local, mas sem ganho de \emph{recall}.
\end{itemize}

\begin{table}[h]
\centering
\caption{Precision / Recall / F1 do detector contra o \emph{ground-truth} sintético
($n=50$ \emph{price shocks}, \texttt{price\_shock}$=0{,}15$).}
\label{tab:metrics}
\begin{tabular}{llccc}
\toprule
\textbf{Ticker} & \textbf{Limiar} & \textbf{Precision} & \textbf{Recall} & \textbf{F1} \\
\midrule
PETR4.SA & estático & 0{,}139 & 0{,}006 & 0{,}011 \\
PETR4.SA & dinâmico & 0{,}548 & 0{,}020 & 0{,}038 \\
VALE3.SA & estático & 0{,}111 & 0{,}005 & 0{,}009 \\
VALE3.SA & dinâmico & 0{,}964 & 0{,}061 & 0{,}115 \\
AMER3.SA & estático & 0{,}694 & 0{,}303 & 0{,}421 \\
AMER3.SA & dinâmico & 0{,}605 & 0{,}120 & 0{,}200 \\
ITUB4.SA & estático & 0{,}769 & 0{,}130 & 0{,}222 \\
ITUB4.SA & dinâmico & 0{,}939 & 0{,}036 & 0{,}069 \\
\bottomrule
\end{tabular}
\end{table}

\subsection{Leitura crítica}
O \emph{recall} baixo \textbf{não} indica falha do detector, mas que a magnitude de injeção é
\textbf{provisória e absoluta}: a dificuldade da tarefa varia arbitrariamente entre ativos,
penalizando os de baixa volatilidade. Esta é a evidência empírica que motiva a recomendação
do ADR-0006: calibrar o choque em $k$ desvios-padrão do retorno \emph{de cada ativo}, tornando
a dificuldade comparável entre setores. A avaliação corrente, portanto, é um piso
metodológico --- valida o \emph{pipeline} de métricas (injeção $\to$ rotulagem $\to$
P/R/F1), e não um resultado final. Adicionalmente, a taxa-base de $\approx$5\% imposta pelo
p95 (Seção~\ref{sec:m4}) limita a \emph{precision} máxima atingível e deve ser lida em
conjunto. A injeção de \emph{volume spikes} (chave \texttt{volume\_spike} no
\texttt{config.yaml}) permanece prevista mas ainda não implementada nesta etapa.

\section{M6 --- Correlação com Eventos e Comparação Setorial}
\label{sec:m6}

\subsection{Validação narrativa}
Complementando a avaliação quantitativa (Seção~\ref{sec:m5}), as anomalias detectadas no
teste (limiar estático p95) são cruzadas com uma linha do tempo curada de eventos
econômicos e políticos brasileiros de 2010--2024 (\texttt{src/events.py}; ADR-0008). Cada
evento é classificado como \emph{sistêmico} (afeta todos os ativos) ou \emph{específico} de
um \emph{ticker}. Considera-se um evento ``coberto'' quando há ao menos uma anomalia dentro
de uma tolerância de $\pm$\texttt{window\_size} ($\pm$30 dias), justificada pelo fato de a
unidade de detecção ser a janela de 30 passos.

A cobertura narrativa global é de $\approx$50\%, mas o número agregado esconde o resultado
relevante --- o \textbf{casamento entre evento e ativo afetado}:
\begin{itemize}[noitemsep]
  \item \textbf{AMER3} cobre 9 de 10 eventos a ela associados, incluindo a fraude contábil
        (jan/2023) e a recuperação judicial --- o sinal mais forte do \emph{dataset}.
  \item PETR4, VALE3 e ITUB4 cobrem cerca de 2 eventos cada (de 6--8): seus modelos disparam
        sobretudo em eventos \emph{sistêmicos} e nos choques setoriais de maior magnitude,
        permanecendo silenciosos em eventos de impacto difuso --- coerente com um detector
        calibrado para a cauda do erro (p95).
\end{itemize}

\begin{figure}[h]
\centering
\includegraphics[width=\textwidth]{m6_amer3_detalhe}
\caption{AMER3 (2022--2024): erro de reconstrução, janelas anômalas (pontos vermelhos) e
eventos (linhas pontilhadas). O pico de erro de jan/2023 coincide com a fraude contábil e o
pedido de recuperação judicial.}
\label{fig:m6-amer3}
\end{figure}

\subsection{Comparação setorial}
Com um modelo por ativo, comparam-se as taxas de marcação entre setores
(Tabela~\ref{tab:setorial}). As taxas reproduzem o achado da M4: PETR4 (energia) e VALE3
(mineração) são estáveis ($\approx$3\%), enquanto AMER3 (varejo) exibe regime de
volatilidade prolongada (26{,}7\% no estático, recalibrado para 14{,}3\% no dinâmico) e
ITUB4 (financeiro) fica em posição intermediária. A diferença reflete o \emph{ativo}, não
necessariamente o \emph{setor} ($n=1$ por setor): trata-se de estudo de caso qualitativo,
não de inferência estatística (ADR-0006).

\begin{table}[h]
\centering
\caption{Taxa de janelas anômalas no teste (2020--2024) por ativo/setor.}
\label{tab:setorial}
\begin{tabular}{llccc}
\toprule
\textbf{Ticker} & \textbf{Setor} & \textbf{Limiar p95} & \textbf{Estático (\%)} & \textbf{Dinâmico (\%)} \\
\midrule
PETR4.SA & Energia/commodities & 0{,}124 & 2{,}9  & 2{,}9  \\
VALE3.SA & Mineração           & 0{,}082 & 2{,}7  & 3{,}4  \\
AMER3.SA & Varejo              & 0{,}077 & 26{,}7 & 14{,}3 \\
ITUB4.SA & Financeiro          & 0{,}088 & 11{,}0 & 4{,}9  \\
\bottomrule
\end{tabular}
\end{table}

\subsection{Contágio sistêmico: o teste da COVID}
Um evento sistêmico deve aparecer em \emph{todos} os ativos. No \emph{crash} da COVID
(fev--mai/2020), os quatro modelos sinalizam anomalia de forma concentrada no período
(PETR4: 35; VALE3: 33; ITUB4: 54; AMER3: 56 janelas), confirmando o contágio entre setores.
Em contraste, eventos setor-específicos --- Brumadinho ($\to$ VALE3), fraude Americanas
($\to$ AMER3), intervenção na Petrobras ($\to$ PETR4) --- isolam-se nos ativos afetados.
A combinação dos dois comportamentos evidencia que os modelos, treinados por ativo, captam
tanto choques de mercado amplos quanto eventos idiossincráticos --- o objetivo de
generalização cruzada declarado na proposta.

\begin{figure}[h]
\centering
\includegraphics[width=\textwidth]{m6_covid_contagio}
\caption{\emph{Crash} da COVID (fev--mai/2020): os quatro ativos, de setores distintos,
sinalizam anomalias concentradas no mesmo período, confirmando o contágio sistêmico.}
\label{fig:m6-covid}
\end{figure}

\paragraph{Limitações em aberto.} O estudo de sensibilidade da janela dinâmica
(\texttt{dynamic\_window} $\in \{126, 252, 504\}$) previsto no ADR-0005 e a calibração das
magnitudes de injeção relativas a $\sigma$ (ADR-0006) permanecem pendentes e são os próximos
refinamentos antes do fechamento do artigo.

\section{Conclusões Parciais e Próximos Passos}
\label{sec:conclusoes}

Até este ponto, o projeto estabeleceu uma base reprodutível (M0) e validou a
qualidade dos dados (M1). Os principais resultados parciais são:
\begin{itemize}[noitemsep]
  \item Infraestrutura com configuração centralizada, \emph{seeds} fixas e proveniência
        documentada de cada hiperparâmetro, separando escolhas fundamentadas de valores
        provisórios.
  \item Quatro séries íntegras (3724 pregões, sem NaN), com anomalias reais já mapeadas
        como âncoras para a validação narrativa.
  \item Esclarecimento do caso AMER3, com decisão registrada de correção do artefato de
        grupamento na próxima etapa.
\end{itemize}

\paragraph{Próximos passos (M2 --- Pré-processamento).} Implementar o cálculo de
log-retornos, a correção do fator de grupamento do AMER3, o \emph{split} temporal com
ajuste do \texttt{MinMaxScaler} restrito ao treino e a geração das janelas deslizantes.
As etapas subsequentes cobrirão o treino do modelo (M3), os limiares de detecção (M4), a
avaliação quantitativa por injeção sintética (M5) e a correlação com eventos e a
comparação setorial (M6).

\paragraph{Registro de decisões.} Todas as decisões de configuração estão documentadas
como ADRs em \texttt{docs/adr/}. Este relatório é atualizado a cada fechamento de
\emph{milestone}, consolidando decisões, fontes, observações e conclusões para compor o
artigo final.


\bibliographystyle{plainnat}
\bibliography{references}

\end{document}
